% =====================================================================================
% sections/appendix_v1.tex: full proof of Theorem thm:v1 (V1 soundness), IEEEtran compsoc
%
% Sources used (paths relative to session 4aa2c028's scratchpad):
%   sp2027_research/v1_design/V1_SPEC.md  Sec. 2 (notation, windows, cut rules, leaves, table,
%     instance layout, the fractional-sum GKR of Sec. 2.4.1, the checks of Sec. 2.5) and Sec. 3
%     (Theorem 1, Lemma 0, Lemma 1, proof steps 1-7, the remarks)
%   wt_sp2027/docs/improvements/C2_v1_proof_size.md  Secs. 3-4 and 9 (checks as implemented,
%     the soundness steps, the review's fixes: P1 clamp within int8, P4 in Kpre, fold aliases,
%     the input range check, the exception bin outside P5; P4 in mode C)
%   sections/defence.tex  (Theorem thm:sound, its notation J, k_j, n_j, t_j and the product form
%     of the column term; the Kpre reuse bound; the Fiat-Shamir argument of sec:fs)
% Labels referenced: thm:v1, lem:v1-logup, sec:v1 (sections/v1.tex); thm:sound, sec:fs,
% sec:setup-proof (sections/defence.tex).
% Labels defined here: app:v1, lem:v1-window, lem:v1-nowrap, lem:v1-gkr, lem:v1-final.
% Uses the macros of sections/v1.tex (re-provided below so that this file also compiles alone).
% =====================================================================================

\providecommand{\Kpre}{\mathsf{K}_{\mathrm{pre}}}
\providecommand{\sC}{\mathsf{C}}
\providecommand{\T}{\top}
\providecommand{\F}{\mathbb{F}_p}
\providecommand{\enc}{\mathrm{Enc}}
\providecommand{\para}[1]{\par\noindent\textbf{#1}\quad}
\providecommand{\qed}{\hfill$\square$}
\providecommand{\Fext}{\mathbb{K}}
\providecommand{\Vlo}{\mathrm{lo}}
\providecommand{\Vshift}{\mathrm{shift}}
\providecommand{\Vrq}{\mathrm{rq}}
\providecommand{\Vtab}{\mathcal{T}}
\providecommand{\Vcut}{\mathcal{C}}
\providecommand{\Vclr}{\mathcal{O}}
\providecommand{\Nleaf}{N_{\mathrm{leaf}}}
\providecommand{\Veq}{\mathrm{eq}}
\ifcsname lemma\endcsname\else\newtheorem{lemma}{Lemma}\fi

\section{Soundness of V1}
\label{app:v1}

This appendix proves Theorem~\ref{thm:v1}. It fixes the notation, states the verifier's checks
and the GKR protocol precisely, proves four lemmas, and then proves the theorem.

\para{Notation.} $p = 15 \cdot 2^{27} + 1$, $\Fext = \F[\xi]/(\xi^8 - 11)$, and $S = \F + \xi\F$ is
the set of elements of $\Fext$ whose coordinates of degree 2 to 7 vanish, so $|S| = p^2$. Integers
are mapped to $\F$ modulo $p$. For a weight operation $\ell$, $A_\ell$ is an integer matrix with
$N_\ell$ rows and $\hat X_\ell = [X_\ell; \mathbf 1]$ is the input that the verifier derives. For
an entry $(i,j)$ of a cut operation, with $i < N_\ell$ and $j < T$, $s_{ij}$ is the sent value,
$L_{ij}$ and $W_{ij}$ are the start and width of its window ($\Vlo_\ell(s_{ij})$ and
$W_\ell(s_{ij})$ for a normal entry, $z$ and 1 for an exception with value $z$), $\delta_{ij} =
(A_\ell \hat X_\ell)_{ij} - L_{ij} \in \mathbb Z$ is its residue and $v_{ij} = \delta_{ij} +
\xi\,W_{ij} \in \Fext$ its leaf. The table is $\Vtab = \{t + \xi w : w \in \mathcal W, 0 \le t < w\}$
and $\mu$ holds the multiplicities. An instance $\gamma$ stacks the rows of its operations,
operation $\ell$ at row offset $o_\ell$, in $R_\gamma$ rows; it has $n_\gamma = n_{\mathrm{row}} +
n_{\mathrm{col}}$ variables with $n_{\mathrm{row}} = \lceil \log_2 R_\gamma \rceil$ and
$n_{\mathrm{col}} = \lceil \log_2 T \rceil$. Leaf index $\iota = \mathit{row} \cdot
2^{n_{\mathrm{col}}} + \mathit{col}$ is \emph{real} if $\mathit{row} < R_\gamma$ and $\mathit{col}
< T$. Bit vectors are indexed from the least significant bit, $\Veq(a, b) = \prod_i (a_i b_i +
(1 - a_i)(1 - b_i))$, and the multilinear extension of a function $g$ on $\{0,1\}^n$ is
$\tilde g(\rho) = \sum_{y} \Veq(\rho, y)\, g(y)$. For an integer $N$, $I_N(r) = \sum_{y < N}
\Veq(r, y)$, which the verifier computes from the binary digits of $N$ in $O(\log N)$ operations.

\para{The verifier's checks.} The verifier rejects at the first failed check.
(D1)~Ranges: $a \in [-127, 127]$, $\Delta \in R_\Delta = [\Vshift_\ell(1 - 2^{29}),
\Vshift_\ell(2^{29} - 1)]$, $|z| < 2^{29}$ for clear claims, and a bound on the inputs of cut
operations (defence in depth only).
(D2)~Derivation: every operation without weights is recomputed; the consumer of a cut operation
reads $s_\ell$.
(D3)~The token rule on the claims of the LM head.
(D4)~Exceptions: each listed entry belongs to a cut requantisation, the indices increase strictly
within an operation, $|z| < 2^{29}$ and $\Vrq_\ell(z) = a_{ij}$.
(D5)~Multiplicities: $|\Vtab|$ entries, $\mu_\tau < p$ for each, and $\sum_\tau \mu_\tau =
\Nleaf$ in the integers.
(G)~Every round check of every instance's GKR protocol, and $\mathsf d_0^\gamma \neq 0$.
(F1)~Freivalds' equation for the clear operations.
(F2)~In $\Kpre$: $\chi_s^\T y_\ell = u_s \bar x_\ell$ for every cut operation and secret vector.
(F3)~For every instance, $\hat{\mathsf n} = I_\gamma$ and $\hat{\mathsf d} = \alpha I_\gamma + 1 -
I_\gamma - \tilde V_\gamma$, where $I_\gamma = I_{R_\gamma}(r_{\mathrm{row}})\,
I_T(r_{\mathrm{col}})$, $\tilde V_\gamma = \sum_{\ell \in \gamma} \big(S_\ell - \chi_\ell^\T
L_\ell e + \xi\, \chi_\ell^\T W_\ell e\big)$, and $S_\ell$ is $u_\ell \bar x_\ell$ for a cut
operation in the row layout and $\chi_\ell^\T y_\ell$ for one in the transposed layout or in
$\Kpre$.
(F4)~$\sum_\gamma \mathsf n_0^\gamma / \mathsf d_0^\gamma = \sum_\tau \mu_\tau / (\alpha -
\tau)$, compared by cross-multiplication.
(C)~The column check of every committed matrix, with the eight coordinate planes of $\chi_\ell$
and $u_\ell$ (cut, row layout), of $\bar x$ and $y$ (cut, transposed layout), or the operands of
v0 (clear).

\para{The GKR protocol of one instance.} Layer $n = n_\gamma$ holds the leaves, $(\mathsf
n_\iota, \mathsf d_\iota) = (1, \alpha - v_\iota)$ for real $\iota$ and $(0, 1)$ otherwise. Node
$y$ of layer $k < n$ combines the children $2y$ and $2y+1$:
\[
  \begin{aligned}
    \mathsf n_k(y) &= \mathsf n_{k+1}(2y)\,\mathsf d_{k+1}(2y{+}1) + \mathsf n_{k+1}(2y{+}1)\,
      \mathsf d_{k+1}(2y), \\
    \mathsf d_k(y) &= \mathsf d_{k+1}(2y)\,\mathsf d_{k+1}(2y{+}1).
  \end{aligned}
\]
(i)~The prover sends $\mathsf n_1(0), \mathsf n_1(1), \mathsf d_1(0), \mathsf d_1(1)$. The
verifier computes the root $(\mathsf n_0, \mathsf d_0)$ by the same rule, draws $\theta_0$, and
sets the claims $\bar{\mathsf n} = \mathsf n_1(0) + \theta_0(\mathsf n_1(1) - \mathsf n_1(0))$
and $\bar{\mathsf d}$ likewise, at the point $\rho = (\theta_0)$.
(ii)~For $k = 1, \dots, n-1$ the verifier draws $\eta_k$. The claim $c = \bar{\mathsf n} +
\eta_k \bar{\mathsf d}$ should equal $\sum_{y \in \{0,1\}^k} \Veq(\rho, y)\, F_k(y)$, where
$F_k(y) = \mathsf n_{k+1}(2y)\mathsf d_{k+1}(2y{+}1) + \mathsf n_{k+1}(2y{+}1)\mathsf d_{k+1}(2y)
+ \eta_k \mathsf d_{k+1}(2y)\mathsf d_{k+1}(2y{+}1)$. A sum-check over the $k$ bits of $y$
follows. In round $j$ the prover sends $g_j(0), g_j(2), g_j(3)$ of a polynomial of degree at most
3; the verifier sets $g_j(1) = c - g_j(0)$, draws $r_{k,j}$, and replaces $c$ by $g_j(r_{k,j})$.
With $r' = (r_{k,0}, \dots, r_{k,k-1})$, the prover then sends $\tilde{\mathsf n}_{k+1}(b, r')$
and $\tilde{\mathsf d}_{k+1}(b, r')$ for $b \in \{0, 1\}$, where $b$ is the lowest bit, and the
verifier checks $\Veq(\rho, r') \big[\tilde{\mathsf n}_{k+1}(0,r')\tilde{\mathsf d}_{k+1}(1,r') +
\tilde{\mathsf n}_{k+1}(1,r')\tilde{\mathsf d}_{k+1}(0,r') + \eta_k \tilde{\mathsf
d}_{k+1}(0,r')\tilde{\mathsf d}_{k+1}(1,r')\big] = c$. It draws $\theta_k$, sets $\bar{\mathsf
n} = \tilde{\mathsf n}_{k+1}(0,r') + \theta_k\big(\tilde{\mathsf n}_{k+1}(1,r') - \tilde{\mathsf
n}_{k+1}(0,r')\big)$ and $\bar{\mathsf d}$ likewise, and moves to the point $\rho \leftarrow
(\theta_k, r')$.
(iii)~The output is the root, the point $\rho = (r_{\mathrm{col}}, r_{\mathrm{row}})$, whose
$n_{\mathrm{col}}$ lowest coordinates form $r_{\mathrm{col}}$, and the final claims $\hat{\mathsf
n} = \bar{\mathsf n}$, $\hat{\mathsf d} = \bar{\mathsf d}$. The transcript has $4 + \sum_{k=1}^{n-1}
(3k+4)$ elements of $\Fext$.

\para{Proof of Lemma~\ref{lem:v1-logup}.} Let $c_v$ be the number of indices $i$ with $v_i = v$,
so that $\sum_i 1/(X - v_i) = \sum_v c_v/(X - v)$, where $c_v$ is read in $\F$. The functions $1/(X
- v)$ for distinct $v$ are linearly independent over $\Fext$ (partial fractions are unique), so the
identity gives $c_v \equiv 0 \pmod p$ for $v \notin \Vtab$ and $c_\tau \equiv \mu_\tau \pmod p$ for
$\tau \in \Vtab$. Since $c_\tau \ge 0$ and $0 \le \mu_\tau < p$, the second congruence gives $c_\tau
\ge \mu_\tau$. Hence $\sum_{\tau \in \Vtab} c_\tau \ge \sum_\tau \mu_\tau = N = \sum_v c_v$, and
$c_v = 0$ for every $v \notin \Vtab$. \qed

\begin{lemma}[windows]
  \label{lem:v1-window}
  For an integer $m_\ell > 0$, $\Vshift_\ell(z) = v$ if and only if $\Vlo_\ell(v) \le z <
  \Vlo_\ell(v+1)$, and $W_\ell(v) \in \{\lfloor 2^{30}/m_\ell \rfloor, \lceil 2^{30}/m_\ell
  \rceil\}$. If (D1) and (D4) pass and every residue of a cut operation $\ell$ lies in $[0,
  W_{ij})$, then $s_\ell$ equals $\Vrq_\ell(A_\ell \hat X_\ell)$ for a requantisation and
  $\Vshift_\ell(A_\ell \hat X_\ell)$ for a residual operation.
\end{lemma}

\noindent\emph{Proof.} $\Vshift_\ell(z) = v$ means $v\,2^{30} \le m_\ell z + 2^{29} < (v+1)\,
2^{30}$, that is, $(v\,2^{30} - 2^{29})/m_\ell \le z < ((v+1)\,2^{30} - 2^{29})/m_\ell$. For an
integer $z$ this is $\Vlo_\ell(v) \le z < \Vlo_\ell(v+1)$. The two bounds differ by
$2^{30}/m_\ell$, so their ceilings differ by $\lfloor 2^{30}/m_\ell \rfloor$ or $\lceil
2^{30}/m_\ell \rceil$. Let $Z = (A_\ell \hat X_\ell)_{ij}$. For a normal entry, $\delta_{ij} \in
[0, W_{ij})$ gives $\Vshift_\ell(Z) = s_{ij}$; for a residual operation this is the claim, and for
a requantisation $|s_{ij}| \le 127$ by (D1), so $\Vrq_\ell(Z) = s_{ij}$. For an exception,
$W_{ij} = 1$ and $\delta_{ij} = 0$ give $Z = z$, and (D4) checked $\Vrq_\ell(z) = a_{ij}$. \qed

\begin{lemma}[no wrap-around]
  \label{lem:v1-nowrap}
  Let $\ell \in \Vcut$ have the honest input, and let (D1) and (D4) pass. If $\delta_{ij} \notin
  [0, W_{ij})$, then $v_{ij} \notin \Vtab$.
\end{lemma}

\noindent\emph{Proof.} By (P4), $|(A_\ell \hat X_\ell)_{ij}| < 2^{29}$. We show $|L_{ij}| < 2^{29}
+ 2^{16}$. For a normal entry of a requantisation, $|s_{ij}| \le 127$ (by (D1), consistent with the
clamp that (P1) requires) and $2^{30}/m_\ell \le 2^{16}$ by (P3), so $|\Vlo_\ell(s_{ij})| \le
127.5 \cdot 2^{16} + 1 < 2^{23}$. For a residual operation, $\Vlo_\ell$ is non-decreasing and
$\Vlo_\ell(\Vshift_\ell(z)) \le z < \Vlo_\ell(\Vshift_\ell(z)) + W_\ell(\Vshift_\ell(z))$ by
Lemma~\ref{lem:v1-window}, so $s_{ij} \in R_\Delta$ gives $1 - 2^{29} - 2^{16} <
\Vlo_\ell(s_{ij}) \le 2^{29} - 1$. For an exception, $|z| < 2^{29}$ by (D4). Hence $|\delta_{ij}|
< 2^{30} + 2^{16}$. If $\delta_{ij} < 0$, then $\delta_{ij} \bmod p = p + \delta_{ij} > p - 2^{30}
- 2^{16} = 939{,}458{,}561 > 2^{16} \ge W_{ij}$. If $\delta_{ij} \ge W_{ij}$, then $\delta_{ij} <
2^{30} + 2^{16} < p$, so $\delta_{ij} \bmod p = \delta_{ij} \ge W_{ij}$. In both cases
$\delta_{ij} \bmod p \notin [0, W_{ij})$. Since $1$ and $\xi$ are linearly independent over $\F$, an
element $t + xw$ of $\Vtab$ equals $v_{ij}$ only if $w = W_{ij}$ and $t = \delta_{ij} \bmod p$,
which would require $\delta_{ij} \bmod p < W_{ij}$. \qed

\begin{lemma}[fractional-sum GKR]
  \label{lem:v1-gkr}
  Fix $\alpha$ and the true leaves of instance $\gamma$, and let $\mathsf n^\ast_k, \mathsf
  d^\ast_k$ be the true layers. If the root $(\mathsf n_0, \mathsf d_0)$ computed in step~(i)
  differs from $(\mathsf n^\ast_0, \mathsf d^\ast_0)$, then, except with probability
  $\varepsilon_{\mathrm{GKR}}(n_\gamma) = \big(1 + \sum_{k=1}^{n_\gamma - 1} (3k + 2)\big)/p^8$
  over the verifier's challenges, either a check (G) fails or $(\hat{\mathsf n}, \hat{\mathsf d})
  \neq (\tilde{\mathsf n}^\ast_{n_\gamma}(\rho), \tilde{\mathsf d}^\ast_{n_\gamma}(\rho))$.
\end{lemma}

\noindent\emph{Proof.} We show that each step keeps a false claim false except with the stated
probability. In step~(i) the root is a function of the four sent values, so at least one of them
differs from its true value. The sent and the true values define lines in $\theta_0$, one pair for
the numerators and one for the denominators, and at least one pair is distinct; two distinct lines
meet at most once, so the claims at $\theta_0$ are false except with probability $1/p^8$. In layer
$k$, if $(\bar{\mathsf n}, \bar{\mathsf d}) \ne (\tilde{\mathsf n}^\ast_k(\rho), \tilde{\mathsf
d}^\ast_k(\rho))$, then $c$ differs from $\tilde{\mathsf n}^\ast_k(\rho) + \eta_k
\tilde{\mathsf d}^\ast_k(\rho)$ except with probability $1/p^8$, since a nonzero affine function
of $\eta_k$ has at most one root. This true value equals $\sum_y \Veq(\rho, y) F^\ast_k(y)$ by
the definition of the tree and of the multilinear extension. The summand, written with
$\tilde{\mathsf n}^\ast_{k+1}(0, y)$, $\tilde{\mathsf n}^\ast_{k+1}(1, y)$ and the corresponding
denominators, has degree at most 3 in each variable, so by the soundness of the
sum-check~\cite{lfkn1992}, except with probability $3k/p^8$ either a round fails or the final $c$
differs from $\Veq(\rho, r')\, F^\ast_k(r')$, the value of the summand at $r'$ with the true
children. If the four sent child values were the true ones, the verifier's check would compare $c$
with exactly this value and fail; so at least one is false, and as in step~(i) the new claims at
$(\theta_k, r')$ are false except with probability $1/p^8$. Each layer $k$ thus adds $(3k +
2)/p^8$ and step~(i) adds $1/p^8$. After layer $n_\gamma - 1$ the claims refer to the leaf layer,
which gives the statement. \qed

\begin{lemma}[final values]
  \label{lem:v1-final}
  With $e$ and $\chi_\ell$ as in Section~\ref{sec:v1}, the true leaf layer of instance $\gamma$
  satisfies $\tilde{\mathsf n}^\ast(\rho) = I_\gamma$ and $\tilde{\mathsf d}^\ast(\rho) = \alpha
  I_\gamma + 1 - I_\gamma - \sum_{\ell \in \gamma} \chi_\ell^\T (A_\ell \hat X_\ell - L_\ell +
  \xi\, W_\ell)\, e$, where $I_\gamma = I_{R_\gamma}(r_{\mathrm{row}})\, I_T(r_{\mathrm{col}})$.
\end{lemma}

\noindent\emph{Proof.} The numerators are 1 on real leaves and 0 elsewhere, so $\tilde{\mathsf
n}^\ast(\rho) = \sum_{\iota \text{ real}} \Veq(\rho, \iota)$. Since $\Veq$ factorises over the
column and row bits and the real leaves form the product of the first $T$ columns and the first
$R_\gamma$ rows, this sum is $I_T(r_{\mathrm{col}})\, I_{R_\gamma}(r_{\mathrm{row}})$. The
denominators are $\alpha - v_\iota$ on real leaves and 1 elsewhere, and $\sum_\iota \Veq(\rho,
\iota) = 1$, so $\tilde{\mathsf d}^\ast(\rho) = \alpha I_\gamma - \sum_{\iota \text{ real}}
\Veq(\rho, \iota)\, v_\iota + (1 - I_\gamma)$. The real leaf in row $o_\ell + i$ and column $j$ is
$v_{ij} = (A_\ell \hat X_\ell - L_\ell + \xi W_\ell)_{ij}$ in $\Fext$, and its weight is
$\Veq(r_{\mathrm{row}}, o_\ell + i)\, \Veq(r_{\mathrm{col}}, j) = (\chi_\ell)_i\, e_j$, which
gives the sum. \qed

\para{Proof of Theorem~\ref{thm:v1}.} We bound the probability that every check passes, so we may
assume that the deterministic checks (D1) to (D5) pass. We follow the order of the checks.

\emph{Step 1: the first wrong operation.} Call a clear operation locally correct if its claims
equal $A_\ell \hat X_\ell$, and a cut operation locally correct if every residue lies in its
window, where $\hat X_\ell$ is the input that the verifier derives. If every operation is locally
correct, then by Lemma~\ref{lem:v1-window} each consumer of a cut operation reads exactly the value
of the honest computation on the same input, and by induction over the graph in topological order
every derived tensor equals the honest one. The LM head is clear, so its claims are the honest
logits, $y = M(x)$, and the token rule holds for them. Hence some operation is locally incorrect.
Let $\ell^\ast$ be the first in graph order; its derived input is the honest input.

\emph{Step 2: a clear $\ell^\ast$.} The argument of Theorem~\ref{thm:sound} applies to its
committed matrix $j \in \mathcal J_{\Vclr}$. The range check keeps the wrong claim distinct from
the true value modulo $p$. With the true folded rows, Freivalds' check fails except with
probability $p^{-r}$; with others, the column check fails except with probability
$\prod_{i<t_j}(k_j - 1 - i)/(n_j - i)$, unless the prover finds a hash collision. In $\Kpre$ the
bound is $p^{-r}$.

\emph{Step 3: a cut $\ell^\ast$ gives an off-table leaf.} Some residue of $\ell^\ast$ lies outside
its window, so by Lemma~\ref{lem:v1-nowrap} its leaf is not in $\Vtab$. The multiset of true leaves
of all cut operations is determined by $C_M$, $x$ and the first message, which fixes the derived
inputs, the sent values and the exceptions. The first message also fixes $\mu$. Both are fixed
before $\alpha$ is drawn.

\emph{Step 4: the logUp identity fails at $\alpha$.} Since (D5) passed, Lemma~\ref{lem:v1-logup}
shows that $\Phi(X) = \sum_i 1/(X - v_i) - \sum_\tau \mu_\tau/(X - \tau)$ is a nonzero rational
function. Every leaf and every table element has zero coordinates of degree 2 to 7, so the poles
of $\Phi$ form a set $P \subseteq S$ with $|P| \le \Nleaf + |\Vtab|$, and $\Phi(X) \prod_{w \in P}
(X - w)$ is a nonzero polynomial of degree below $|P|$. The challenge $\alpha$ is uniform on
$\Fext \setminus S$, which has $p^8 - p^2$ elements and contains no pole, so $\Phi(\alpha) = 0$
with probability at most $(\Nleaf + |\Vtab|)/(p^8 - p^2)$. Otherwise $\sum_\gamma \mathsf
n^{\ast\gamma}_0 / \mathsf d^{\ast\gamma}_0 = \sum_i 1/(\alpha - v_i) \ne \sum_\tau \mu_\tau /
(\alpha - \tau)$, where every true root denominator is a product of nonzero factors $\alpha - v_i$.
If (F4) passes, with $\mathsf d_0^\gamma \ne 0$ by (G), some instance $\gamma$ has $\mathsf
n_0^\gamma / \mathsf d_0^\gamma \ne \mathsf n^{\ast\gamma}_0 / \mathsf d^{\ast\gamma}_0$, so its
root is not the true root.

\emph{Step 5: the error reaches the final claims.} By Lemma~\ref{lem:v1-gkr}, except with
probability $\varepsilon_{\mathrm{GKR}}(n_\gamma)$, either (G) fails or the final claims of
$\gamma$ differ from the true values at $\rho$.

\emph{Step 6: the final check.} By Lemma~\ref{lem:v1-final}, if $S_\ell = \chi_\ell^\T A_\ell \hat
X_\ell e$ for every $\ell \in \gamma$, the right-hand sides of (F3) are the true values, and (F3)
rejects. Otherwise some $S_\ell$ is wrong, so $u_\ell \ne \chi_\ell^\T A_\ell$ (row layout) or $y
\ne A \bar x$ (transposed layout, or $\Kpre$). In mode $\sC$ with the row layout, write $\chi_\ell
= \sum_{q < 8} \xi^q \chi^{(q)}$ and $u_\ell = \sum_{q<8} \xi^q u^{(q)}$, where $\chi^{(q)}$ and
$u^{(q)}$ have entries in $\F$; then $u^{(q)} \ne \chi^{(q)\T} A_\ell$ for some $q$. Because $C_M$
is honest, every row of the encoded matrix $\hat A$ is a codeword, so $\chi^{(q)\T} \hat A =
\enc(\chi^{(q)\T} A_\ell)$ and $\enc(u^{(q)})$ are distinct codewords that agree on at most $k_j -
1$ of the $n_j$ positions. The $t_j$ distinct columns are drawn after $u_\ell$ is sent, so they all
fall on agreeing positions with probability at most $\prod_{i<t_j}(k_j - 1 - i)/(n_j - i)$;
otherwise (C) fails, unless the prover finds a hash collision. The transposed layout is symmetric:
$\bar x^{(q)\T} \hat A' = \enc(A \bar x^{(q)})$ differs from $\enc(y^{(q)})$ for some $q$. In
$\Kpre$, $y_\ell - A_\ell \bar x_\ell$ has a nonzero coordinate plane $d \in \F^{N_\ell}$, and
$\chi_s^\T (y_\ell - A_\ell \bar x_\ell) = 0$ requires $\chi_s^\T d = 0$, which holds with
probability $1/p$ for a uniform secret vector $\chi_s$ that the prover has not seen; over the $r$
independent secret vectors this is $p^{-r}$.

\emph{Step 7: union bound.} Adding the probabilities of steps 2 and 4 to 6 over the operations,
the instances and the committed matrices gives the bound of Theorem~\ref{thm:v1}. \qed

\para{Reuse in $\Kpre$.} As in v0, the prover learns each verdict and could repeat an accepted
forgery. Until one is accepted, the verdicts do not depend on the secret vectors, since $y_\ell$,
the GKR transcripts and $\alpha$ do not involve them. A false acceptance among $Q$ queries that
reuse the vectors therefore occurs with probability at most $Q$ times the bound.

\para{Fiat--Shamir.} The interactive protocol is round-by-round sound, as in
Section~\ref{sec:fs}: once a message makes the statement doomed, only the next challenge can rescue
it. For the new rounds this probability is at most $1/p^8$ for $\theta_k$ and $\eta_k$, $3/p^8 =
2^{-245.7}$ for a sum-check round, and $(\Nleaf + |\Vtab|)/(p^8 - p^2) \le 2^{-215.3}$ for
$\alpha$. All are below the per-round error that the Fiat--Shamir parameters of v0 allow at
$\lambda = 128$, so its state-restoration bound for $Q$ oracle queries carries over unchanged.

\para{Remarks.} (i)~\emph{Saturated entries.} If the prover omits an exception, the entry keeps its
normal window, which the true pre-activation misses; if it lists a false value $z$, the residue is
nonzero in a window of width 1. Either way Step~3 applies. (ii)~\emph{Residual operations.}
$\Delta$ is unclamped, so its windows are those of $\Vshift_\ell$; the verifier applies the clamp
of the sum itself in (D2). (iii)~\emph{Input bounds.} The proof does not use the input range check
of (D1), because Step~1 works with the honest input of $\ell^\ast$; the check is defence in depth.
(iv)~\emph{The claim bound (P4).} In $\Kpre$ the verifier evaluates (P4) from its own weights. In
mode $\sC$ it holds no weights, and (P4) is a property of the committed model, checked when the
honest commitment is computed, as is v0's claim range. With an untrusted commitment
(Section~\ref{sec:setup-proof}) the bound must be published with the model. (v)~\emph{Completeness.}
Honest multiplicities are below $2^{29} < p$ by (P5). The exception bin of width 1 is not covered
by (P5), so a proof with more than $2^{29}$ clamped entries cannot be encoded; this affects
completeness only.
